\documentclass{article}
\usepackage{icmlko}
\usepackage{graphicx}
\usepackage{booktabs}
\usepackage{amsmath}

\title{닿지 않는 편이 나았다 \\
  {\large 새 자료 0 으로 열리는 레버 둘을 당겼더니, 둘 다 아래로 갔다}}

\begin{document}
\maketitle

\begin{abstract}
직전 사이클은 「모형에 안 닿는 축 칸」을 305 개 세고, 그중 \textbf{새 자료 없이 열 수 있는} 레버를
여덟 개로 정리했다. 이 노트는 그중 값이 제일 큰 둘을 실제로 당겼다.
게임 전용 축 넷을 정식화로 닿게 하는 길 셋(\texttt{union} · \texttt{DOMAX} · NaN 패딩)과,
시장팝업 결측 88 행을 채워 \texttt{BLOCK} 을 떼는 길이다.
\claimx{둘 다 반증됐다.}{어느 한 팔이라도 짝 $\Delta$ 의 BCa 하한이 0 위로 올라오거나, \texttt{BLOCK} 해제+채움이 27 이상의 고유 설계행을 내면 이 문장은 무너진다.} 축 넷은 12 씨앗 짝 $\Delta$ 로 세 팔 모두 상한이 0 아래이고
(\,$-0.001231$ · BCa $[-0.001663,-0.000802]$\,), 결측 채움은 유보 설계 해상도를
$27 \rightarrow 21$ 로 \textbf{내린다}. 그리고 손실의 출처가 뜻밖이었다 —
축 넷을 닿게 했을 때 제일 크게 지는 것은 \textbf{그 넷을 가진 게임 자신}이다(손실의 76.5\%).
\end{abstract}

\section{물음: 「닿지 않는다」는 손실인가}

노트 899 는 축 칸 436 개 중 305 개가 모형에 안 닿는다고 셌고, 이번 사이클의 수리 팔이
다섯째 갈래 ⓓ 「정식화가 안 본다」를 넣어 309 로 고쳤다. ⓓ 는 정확히 넷이다 —
\texttt{price} · \texttt{age\_rating} · \texttt{ram\_gb} · \texttt{n\_category} 로 \textbf{전부 게임에만} 있고,
\texttt{common}(36) 과 \texttt{union}(40) 의 차이 그 자체다. 이 넷은 유보에서 \textbf{살아 있다}
(마스크 1.0/1.0/0.9667/1.0). 그래서 899 의 배정 체계는 이들을 「살아 있다」로 세고
죽은 칸 305 에 한 칸도 안 넣었다.

이 구도가 매력적이었던 이유는 분명하다: \textbf{새 자료가 0 이다.} 자료는 이미 저장소에 있고
막는 것은 \texttt{lab/forms.py:197} 의 \texttt{AXIS\_MODE} 한 줄이다.
같은 성질을 가진 둘째 레버가 시장팝업 \texttt{entry\_friction} 이었다 —
\texttt{lab/fixaxes.py:119-120} 의 \texttt{BLOCK} 한 줄이 막고 있고, 결측 88 행만 채우면
그 줄을 뗄 자격이 생긴다.

이 노트의 물음은 하나다. \textbf{안 닿던 것을 닿게 하면 판이 오르는가.}

\section{사전등록과 배선 검사}

사전등록 \texttt{docs/prereg\_900\_levers.md} 는 커밋 \texttt{f3ca1d3ea}
(2026-08-10T16:34:28Z)로 \textbf{측정보다 앞서} 들어갔다. 측정 산출물 스탬프는 16:47:26Z·17:05:08Z 다.

배선 검사(\texttt{runners/wire900.py})는 887 형 중립화 — 축을 주입해도 조용히 무효화되는 사고 —
를 먼저 배제했다. 팔별로 \texttt{fit} 이 실제로 본 설계 열을 셌다:
A(챔피언) 88 열 · B(\texttt{union}) 96 · C(\texttt{DOMAX}) 96 · D(NaN 패딩) 96.
\claimx{게임 4 열은 진짜로 모형에 닿는다 --- 88 에서 96 으로 늘고, 세 팔 모두 정확히 +8 이다.}{팔 B/C/D 의 설계 열이 88 로 남거나, \texttt{fit} 이 센 열과 \texttt{\_design} 이 센 열이 어느 도메인에서든 어긋나면 이 배선 검사는 무효다(887 형 중립화).}
\texttt{fit}(\texttt{forms.py:669-683})이 \texttt{\_design}(\texttt{:1101})을 안 부르고 같은 조립을
다시 적어 놓았으므로 둘을 따로 세서 대조했고 전 팔·전 도메인에서 일치했다.
유보 가중 합 3{,}775 · 12 도메인 · 자료 지문 \texttt{721a7b7e0d15} 는 두 실행에서 같다.

\subsection{「도메인별로 자기 축만 보는 정식화」는 있지만, 그것도 패딩을 못 피한다}

\texttt{BagBoost.DOMAX}(노트 529 · \texttt{forms.py:837})가 그 장치다. 그러나 \texttt{\_domax}
(\texttt{:976})는 모든 도메인에 모든 짝의 열을 만들고 남의 자리를 0.5/0 으로 채운다
(코드 주석: \emph{"적합과 예보가 같은 자리를 써야 하므로"}).
\claimx{실측 증거 --- 팔 B(\texttt{union})와 팔 C(\texttt{DOMAX})의 열 집합이 12 도메인 전부에서 같다.}{어느 한 도메인에서라도 팔 B 와 팔 C 의 열 집합이 달라지면 \texttt{DOMAX} 는 순서 치환이 아니라 진짜 도메인별 정식화다.}
단일 풀링 설계행렬에서 열 수는 도메인마다 다를 수 없다. \texttt{DOMAX} 는 \texttt{union} 의
열 순서 치환이다. 그래서 패딩 값 자체를 의심하는 팔 D(0.5 대신 NaN)를 따로 세웠다.

\section{판정: 세 팔 전부 패}

\begin{table}[h]\centering\small
\begin{tabular}{lrrrl}
\toprule
팔 & 판 짝 $\Delta$ & BCa 95\% & 양수 씨앗 & 규약 47 \\
\midrule
B \texttt{union}  & $-0.001231$ & $[-0.001663,-0.000802]$ & 1/12 & 패 \\
C \texttt{DOMAX}  & $-0.001224$ & $[-0.001657,-0.000792]$ & 1/12 & 패 \\
D NaN 패딩        & $-0.001503$ & $[-0.002589,-0.000493]$ & 3/12 & 패 \\
\bottomrule
\end{tabular}
\end{table}

48 적합 · 조용한 탈락 0 · 전 팔 전 씨앗 12 도메인·가중 3{,}775 온전.
\claimx{상한이 0 아래이므로 이것은 「모른다」가 아니라 패다(규약 47).}{구간이 \texttt{cluster\_boot} BCa 가 아니라 폴백이었거나, 재표집 단위를 행/날짜로 바꿨을 때 상한이 0 위로 올라오면 「패」가 아니라 「모른다」로 내려야 한다.}
구간은 전부 \texttt{lab/pairboot.cluster\_boot} BCa 이고 폴백은 없다.

\begin{figure}[h]\centering
\includegraphics[width=\linewidth]{figs/fig1.pdf}
\caption{(a) 축 넷을 닿게 한 세 방식 모두 팔 A 아래. (b) 짝 $\Delta$ 와 BCa 를
채택 문턱 0.00353 에 댄 그림 --- 셋 다 상한이 0 아래다. (c) 손실의 76.5\% 가
축 넷의 주인인 게임 자신에게서 나온다.}
\end{figure}

\subsection{넘으려면 얼마여야 하나}

직전 사이클이 안 쓴 줄을 이번엔 쓴다. 게임 하나로 판 문턱을 넘으려면
\textbf{게임 $\Delta\rho \geq +0.07403$} 이어야 한다(일 군집 문턱이면 $+0.07760$).
실측은 $-0.019781$ 로 필요치의 $-0.267$ 배다.
\claimx{크기가 모자란 게 아니라 부호부터 반대다.}{게임 $\Delta\rho$ 의 BCa 상한이 0 위로 올라오면 이 문장은 「크기가 모자라다」로 바뀐다. 지금은 상한 $-0.015442$ 로 0 에 닿지 않는다.}
$|$판 $\Delta|$ 는 채택 문턱의 0.35 배다 --- 씨앗 짝 자로는 확실한 음수이고,
채택 문턱으로는 크기가 그 아래다. 두 문장은 서로 다른 자의 것이므로 같이 적는다.

\subsection{제일 크게 지는 것은 게임 자신이다}

\claimx{「0.5 패딩이 남의 도메인을 해친다」는 그림은 안 나왔다.}{나머지 열하나의 판 기여 합이 게임 기여와 같은 크기로 커지거나, 음수 도메인이 11/11 로 몰리면 패딩 가설이 되살아난다.}
게임 $\Delta\rho = -0.019781$ (BCa $[-0.024225,-0.015442]$ · \textbf{0/12 양수} — 모든 씨앗에서 내려간다)이고
판 기여 $-0.000943$ 으로 손실의 76.5\% 다. 나머지 열하나의 판 기여 합은 $-0.000289$ 로 작고,
그중 $\Delta\rho$ 가 음수인 것은 6/11 이며 아이돌은 오히려 $+0.0081$ 이다.
사전등록 5.3 이 미리 적어 둔 두 갈래 중
\textbf{「축을 닿게 해도 그 도메인조차 안 산다」} 쪽이 걸렸다.

\claimx{팔 D 가 팔 B 보다 더 나쁘다 --- 게임엔 NaN 칸이 하나도 없는데도.}{게임 행에서 NaN 칸이 하나라도 세어지면(지금 0 개) 이 추론은 무너진다 --- 그때는 게임 자신의 NaN 이 원인일 수 있다.}
게임 $\Delta\rho$ 가 $-0.0198 \rightarrow -0.0299$ 로 벌어진다. 게임 행에는 NaN 이 0 개이므로
이 차이는 \textbf{남의 도메인 패딩 값이 게임 행의 예보를 바꾼다}는 뜻이다.
899 의 「축이 있는 척한다」를 NaN 으로 정직하게 고쳐도 안 낫는다.

\section{둘째 레버: 빈칸은 구멍이 아니라 표지였다}

\begin{figure}[h]\centering
\includegraphics[width=\linewidth]{figs/fig2.pdf}
\caption{(a) 유보 126 행의 고유 설계행. 27 은 수집이 아니라 \texttt{BLOCK} 한 줄에서 나오고,
채우면 21$\sim$23 으로 내려간다. (b) 필요치와 실측.}
\end{figure}

시장팝업 결측을 전수로 셌다. 조항 59 를 먼저 걸어 「없다」의 가짓수를 확인했다 —
647 파일에서 \textbf{키 없음 0 · 빈 문자열 0}, 전부 \texttt{null} 한 가지다.
그리고 세 가지가 정정됐다: 결측 37 행은 MKT2 만이 아니라 MKT 15 + MKT2 22 이고,
MKT2 는 37 행이 아니라 239 파일/44 축행/22 유보이며,
\claimx{채워야 할 것은 37 이 아니라 88 이다}{학습 51 행의 결측이 채점에 안 닿는다는 것이 보여지면 37 로 충분하다. 지금은 마스크가 학습에서도 라벨과 $+0.1983$ 으로 이어져 있다.} --- 학습 51 을 두면 표지가 학습/유보로 옮겨갈 뿐이다.

원인은 하나가 아니라 둘이다. MKT2 22 행은 \texttt{ingest/market\_merge.py:150-152} 가
\texttt{None} 을 상수로 박고 2 차 원천에 \texttt{intervention} 블록이 아예 0 건이라
\textbf{파서가 흘린 게 아니라 프롬프트가 안 물었다.} MKT 15 행은 카테고리로 갈린다
(fashion 59/82 = 72.0\% 대 beauty 3/39 = 7.7\%).

\subsection{채우면 내려간다}

챔피언 씨앗 0 을 재적합해 유보 126 행의 고유 설계행을 셌다:
현행 15 · \texttt{BLOCK} 만 항등 \textbf{27}(새 자료 0) · 해제+채움 \textbf{21}(씨앗 0) / 23(씨앗 7) ·
채움만(음성대조) 15.
\claimx{27 은 수집이 아니라 \texttt{fixaxes.py:119-120} 한 줄을 떼면 그 자리에서 나온다.}{\texttt{BLOCK} 만 항등으로 바꾼 변형(새 자료 0)이 15 에 머물면 이 문장은 거짓이다.}
\claimx{그리고 채우면 오히려 21 로 내려간다.}{다른 씨앗이나 다른 채움 값 아래에서 27 이상이 나오면 무너진다. 지금 본 것은 씨앗 둘(21 · 23)뿐이라 이 반증 조건은 아직 열려 있다.}
\texttt{\_feat} 무늬가 안 채우면 넷 $(0.35,1)(0.8188,1)(0.9688,1)(0.5,0)$, 채우면 셋이고,
사라지는 넷째 $(0.5,0)$ 이 \textbf{「이 행은 안 물어본 배치다」라는 결측 표지 그 자체}다.
$27-21=6$ 은 그 패턴이 막으려던 누출에서 나온다.

그렇다고 \texttt{BLOCK} 을 그냥 뗄 수는 없다: 배치$\leftrightarrow$마스크 $-0.6266$
(유보 $-0.7133$)인데 마스크$\leftrightarrow$라벨이 학습에서 $+0.1983$ 이다.
\claimx{채움의 값어치는 설계 패턴 수가 아니라 「\texttt{BLOCK} 을 뗄 자격」에 있다.}{마스크와 라벨의 학습 상관이 0 근처로 내려가면 \texttt{BLOCK} 은 그냥 떼도 된다 --- 그러면 채움은 값어치가 없다.}

\section{자가 적발}

\textbf{하나 --- 사전등록이 죽은 숫자를 인용했다.} §1 이 판 정본을 \textbf{0.46982}
(\texttt{out112\_board.json})로 적었는데 그건 은퇴한 「서수 순위 시대 값」이다.
노트 898(커밋 \texttt{39afa03e6})이 \texttt{state/rank\_test.spearman} 을 서수에서 동률 평균으로
바꾸며 0.46982 를 0.47034 로 갈아 끼웠다. 판정 영향은 없다(짝 $\Delta$ 는 같은 실행·같은 자료).
그 필드가 낸 차 $+0.00051922$ 는 898 이 등록한 「스피어만 구현」 성분과 같은 수다.
\claimx{저장소 카드 본문은 이미 옳았고, 옛 값이 남아 있던 자리는 새 세션이 제일 먼저 읽는
색인 한 줄이었다 --- 그래서 그대로 베껴졌다.}{커밋 시각으로 되짚어 카드 본문이 은퇴 시점 이후에도 옛 값을 이고 있었다면 원인은 색인 한 줄이 아니라 카드 자체다.}

\textbf{둘 --- 결측 채움 시험이 한 번도 안 걸렸다.} 1 차 시험의 채움 패치가 \texttt{json.load} 를
걸었는데 \texttt{state/tri\_domain.py:69} 는 \texttt{json.loads(read\_text())} 라 무연산이었다.
마스크 합이 89 로 안 변한 것으로 잡았고, 재시험에 \texttt{assert 채운 행 == 88} 을 박았다.
\claimx{실패한 시험도 증거물이라 두 파일 다 남겼다}{1 차 산출물이 지워지거나 덮이면 「27 이 무효값」이라는 사실을 뒤에서 확인할 길이 없다.} --- 1 차의 「27」은 채움이 안 걸린 무효값이다.

\textbf{셋 --- 게이트를 안 돌린 채 「고쳤다」고 적을 뻔했다.} \texttt{ingest.audit.dead\_numbers()}
를 돌리니 이 사이클의 파일이 경성 4 건으로 걸렸다.

\section{한계 --- 모른다고 적는 것들}

\begin{itemize}
\item \textbf{게임이 왜 지는지 모른다.} (가) 넷이 라벨과 무관 (나) 패딩 열이 그 넷을 덮음.
      이 사이클은 둘을 못 가른다 --- 사전등록에 미리 적어 둔 문장 그대로다.
\item 학습-유보 차를 안 쟀다(과적합과 분포이동을 못 가른다). 4 축 삭제시험도 안 했다
      (넷을 한 덩어리로만 켰다). 둘 다 사후 팔 금지라 다음 사이클 몫이다.
\item 채운 값이 가정이라 21 대 23 이 흔들린다(씨앗 둘만 봤다).
      채움이 판 $\rho$ 에 주는 영향은 \textbf{안 쟀다}.
\item 바깥 확인은 6 건 표본이지 37 행 실측이 아니다(조항 60).
      \texttt{m.blog.naver.com} 은 이 환경에서 못 가져온다 --- 「없다」가 아니라 「못 읽었다」다.
\item 구간의 재표집 단위는 씨앗 12 다. 노트 897 의 날짜 군집 540 을 재는 자가 아니다.
\item 두 문턱(0.00353 · 0.00370) 모두 오늘도 12 도메인 중 \textbf{0 개}를 통과시킨다.
      이번 팔도 그 자를 못 흔든다.
\end{itemize}

\section{남는 것}

레버 8 개 표에서 게임 넷이 지워진다 --- \textbf{자료가 없어서 못 쓴 게 아니라, 써도 판이 내려간다.}
시장팝업 하나는 성격이 바뀐다: 값어치가 설계 패턴이 아니라 \texttt{BLOCK} 해제 자격이므로,
88 행을 채우는 일은 「해상도를 사는 것」이 아니라 「누출 없이 한 줄을 뗄 권리를 사는 것」이다.
남은 것은 펀딩 \texttt{trend}$\times$3 (449 행 · 판 가중 529 로 최대)이다.

\end{document}
